% Part: first-order-logic
% Chapter: introduction
% Section: substitution

\documentclass[../../../include/open-logic-section]{subfiles}

\begin{document}

\olfileid[pt]{fol}{int}{sub}

\olsection{Substituição}

Discutiremos um exemplo para ilustrar como as coisas se articulam e
como o desenvolvimento da sintaxe e da semântica estabelece as bases para
nossas investigações mais avançadas posteriormente. Nossos sistemas de !!{derivation} deveriam permitir-nos
derivar $\Atom{\Obj P}{\Obj a}$ de $\lforall[\Obj
v_0][\Atom{\Obj P}]{\Obj v_0}$. Talvez até queiramos afirmar isso como uma
regra de inferência. No entanto, para fazê-lo, devemos ser capazes de enunciá-lo
nos termos mais gerais: não apenas para $\Obj P$, $\Obj a$ e $\Obj
v_0$, mas para qualquer !!{formula}~$!A$, e termo~$t$, e
!!{variable}~$x$. (Lembre-se de que !!{constant}s são termos, mas vamos
considerar também termos mais complicados construídos a partir de !!{constant}s e
!!{function}s.) Então, queremos ser capazes de dizer algo como,
``sempre que você tiver derivado $\lforall[x][!A(x)]$ você está justificado
ao inferir ~$!A(t)$---o resultado da remoção de $\lforall[x]$ e
substituindo ~$x$ por ~$t$.'' Mas o que exatamente significa ``substituir $x$ por ~$t$''?
Qual é a relação entre $!A(x)$ e ~$!A(t)$?  Isso
sempre funciona?

Para tornar isso preciso, definimos a operação de \emph{substituição}.
A substituição é realmente complicada, porque não podemos simplesmente substituir
todos os ~$x$ em ~$!A$ por ~$t$, e nem todos os ~$t$ podem ser substituídos por
qualquer ~$x$. Lidaremos com isso novamente usando definições indutivas. Mas
uma vez feito isso, especificando uma regra de inferência como ``inferir $!A(t)$
de $\lforall[x][!A(x)]$'' torna-se uma definição precisa.  Além disso,
seremos capazes de mostrar que esta é uma boa regra de inferência no sentido
de que $!A(t)$ é consequência de $\lforall[x][!A(x)]$. Mas para provar isso, temos
de provar novamente algo que à primeira vista pode levá-lo a perguntar
``por que estamos fazendo isso?'' Que $!A(t)$ seja consequência de $\lforall[x][!A(x)]$
depende do fato de que se $\Sat{M}{!A(t)}$ vale ou não depende
apenas do valor do termo ~$t$, ou seja, se deixarmos $m$ ser qualquer
!!{element} de ~$\Domain{M}$ que for selecionado por ~$t$, então
$\Sat{M}{!A(t)}[s]$ se, e somente se, $\Sat{M}{!A(x)}[\Subst{s}{m}{x}]$. Isso vale
mesmo quando $t$ contém !!{variable}s, mas teremos que ter cuidado
com a forma exata como declaramos o resultado.

\end{document}
